\section{A selected Dionysian--Thomistic concordance}
\label{app:dionysian-concordance}

This concordance records passages used in the exposition. ``Express'' means that the \emph{Summa} names Dionysius and the relevant chapter; ``parallel'' marks a comparison of arguments. An objection is identified as an objection and must be read with its reply. The table is a reading guide, not a count of all Dionysian citations in Aquinas. Parker page ranges refer to the 1899 part~II; PG~3 ranges identify the original chapter text, excluding subsequent commentary where possible.

\begingroup
\small
\setlength{\tabcolsep}{4pt}
\begin{longtable}{@{}>{\raggedright\arraybackslash}p{.06\linewidth}>{\raggedright\arraybackslash}p{.17\linewidth}>{\raggedright\arraybackslash}p{.32\linewidth}>{\raggedright\arraybackslash}p{.38\linewidth}@{}}
\toprule
CH & Parker / PG~3 & \emph{Summa theologiae} & Reception or qualification \\
\midrule
\endhead
I & 1--4 / 119--124 & I q.~111 a.~1, body & Express: human illumination through sensible likenesses. \\
II & 4--13 / 135--146 & I q.~1 a.~9 ad~3 & Dissimilar symbols protect against literalism. The received English reference prints CH~I; the developed source argument is II. \\
III & 13--16 / 163--168 & I q.~108 a.~1, objection~2 and body; q.~106 a.~2 ad~1 & Express hierarchy definition; divine Persons do not undergo purification. Angelic purification can remove ignorance without sin. \\
IV & 16--21 / 177--182 & I q.~111 a.~1, sed contra; q.~106 a.~3, sed contra & Express: angelic mediation and orderly illumination. \\
V & 21--22 / 195--196 & I q.~108 a.~5, body and ad~1 & Parallel: common names and proper perfections, distinguished by participation and eminence. \\
VI & 23--24 / 199--202 & I q.~108 a.~1, sed contra; a.~2, body; a.~3, body & Express triads; our general distinctions do not exhaust each angel's particular rank. \\
VII & 24--31 / 205--212 & I q.~108 a.~5 ad~5--6; q.~106 a.~1; q.~107 a.~2, sed contra & Express names, illumination, and questions addressed upward. All blessed angels see God's essence immediately. \\
VIII & 31--35 / 237--242 & I q.~106 a.~1, sed contra; q.~108 a.~5 ad~1--3 & Express communication of light and explanation of the middle orders. \\
IX & 35--40 / 257--262 & I q.~108 a.~5 ad~3--4; q.~113 a.~3 & Express names; comparison with the distinction between corporate and individual guardianship. \\
X & 41--42 / 271--274 & I q.~108 a.~3 ad~1 & Express: first, middle, and last within the same order. \\
XI & 42--44 / 283--286 & I q.~54 a.~3, sed contra; q.~108 a.~5 ad~1 & Express: essence, power, and operation; common and special senses of power. \\
XII & 44--45 / 291--294 & I q.~55 a.~3, sed contra; q.~106 a.~4, objection~1 and reply & Express: more universal knowledge; differing reception without a superior's withholding illumination. \\
XIII & 46--53 / 299--308 & I q.~112 a.~2, body and ad~2 & Express: Isaiah's minister explained by shared office or mediated agency. \\
XIV & 53--54 / 321--322 & I q.~50 a.~3, body; q.~112 a.~4 ad~2 & Express multitude; interpretation of those who stand before God and those who minister. \\
XV & 54--66 / 325--340 & I q.~106 a.~1, body; q.~107 a.~5, objection~3 and reply & Express: teeth as the adaptation of intellectual nourishment; shared illumination and particular speech. \\
\bottomrule
\end{longtable}
\endgroup

CH V and XIV have no numbered subdivisions in Parker. Chapter VI is also printed without section headings there, whereas PG~3 marks sections~1--2. Those differences concern the editions' navigation. They do not supply additional chapters or alter the order of the argument.
